% Chapter 4, Figure 14: non-vertical trajectories for a 30 deg launch angle, F7 engine (scan: figures/ch4/fig14.png).
% Curves: fig14-a/b/c.csv (fig14.py), DIGITIZED from the 1973 art: the book gives no F7 engine model, so its 2-D
% method, eqs. (106)-(115) with (125)-(131), cannot be run (corrections/v2-figures.md). The times (seconds from
% launch) are as printed; their leaders touch the curves at the apex and the impact of each (fig14-marks.csv: the
% highest point and the last point of each digitized curve). Curve (c) strikes the ground under power at 6.84 s
% (caption); its tag stands on the baseline of its apex time, as printed, the curve (2 mm high) being too small
% to carry it. The curves are read in the family's frame convention (fig10.py: x between the drawn x ticks, along
% the drawn y axis; the 1973 frame is slightly skewed and its x scale uneven, lettering and all).
% The trajectory family (Figs 10-14): open axes, equal x and y scales as printed (the axes 4.2 in wide), the printed
% ticks; each axis line runs half a tick step past its last label where a curve does (here curve (a): apex 1354 m,
% impact 2058 m); curves (a), (b), (c) solid in s1, s2, s3, each tag 2.75 mm out along the normal of its
% descending branch (its centre computed by fig14.py, in fig14-marks.csv); times \footnotesize in ink on leaders.
\documentclass[11pt]{standalone}
\usepackage{tamrfig}
\pgfplotstableread[col sep=comma]{figures/v2/ch4/fig14-marks.csv}\trajmarks
\newcommand{\getmark}[3]{\pgfplotstablegetelem{#1}{#2}\of\trajmarks \edef#3{\pgfplotsretval}}
\getmark{0}{apexx}\Aax \getmark{0}{apexy}\Aay \getmark{0}{impactx}\Aix \getmark{0}{tagx}\Atx \getmark{0}{tagy}\Aty
\getmark{1}{apexx}\Bax \getmark{1}{apexy}\Bay \getmark{1}{impactx}\Bix \getmark{1}{tagx}\Btx \getmark{1}{tagy}\Bty
\getmark{2}{apexx}\Cax \getmark{2}{apexy}\Cay \getmark{2}{impactx}\Cix
% \timemark{<x>}{<y>}{<angle>}{<length>}{<label anchor>}{<time>}: a leader from the point on the curve to the time
\newcommand{\timemark}[6]{%
  \draw[leader] (axis cs:#1,#2) -- ++(#3:#4) node[anchor=#5, font=\footnotesize, inner sep=1.2pt] {#6};}
\begin{document}
\begin{tikzpicture}
  \begin{axis}[tamr,
      x={4.2in/2200}, y={4.2in/2200},
      xmin=0, xmax=2200, ymin=0, ymax=1400,
      xtick={0,400,...,2000}, ytick={0,400,...,1200},
      xlabel={$x$ (m)}, ylabel={$y$ (m)}]
    \addplot[series1, solid] table[x=x, y=y, col sep=comma] {figures/v2/ch4/fig14-a.csv};
    \addplot[series2, solid] table[x=x, y=y, col sep=comma] {figures/v2/ch4/fig14-b.csv};
    \addplot[series3, solid] table[x=x, y=y, col sep=comma] {figures/v2/ch4/fig14-c.csv};
    % times from launch at the apex and the impact, as printed
    \timemark{\Aax}{\Aay}{-97}{4.5mm}{north}{15.20}
    \timemark{\Aix}{0}{145}{6mm}{south east}{41.03}
    \timemark{\Bax}{\Bay}{60}{5mm}{south west}{10.40}
    \timemark{\Bix}{0}{60}{6mm}{south west}{20.01}
    \timemark{\Cax}{\Cay}{19}{9.5mm}{south west}{3.80}
    \timemark{\Cix}{0}{10}{6mm}{south west}{6.84}
    \node[curve tag] at (axis cs:\Atx,\Aty) {a};
    \node[curve tag] at (axis cs:\Btx,\Bty) {b};
    % the (c) tag on the 3.80 baseline, as printed (from the label's anchor, 7.66 mm right), 6.84 below it, as
    % printed; the 3.80 leader is long enough to lift the pair clear of 6.84 (0.9 mm), its label clear of curve (b)
    % (1.7 mm), the tag 3.9 mm clear of curve (b)
    \node[curve tag, anchor=south west] at ([shift={(19:9.5mm)}, xshift=7.66mm]axis cs:\Cax,\Cay) {c};
  \end{axis}
\end{tikzpicture}
\end{document}
